% =========================================================
% 6.3 DISEÑO ESTRUCTURAL Y DE COMPORTAMIENTO
% El \section{} de esta sección está en el chapter.tex del capítulo;
% sus referencias van en seccion6_3.bib (misma carpeta).
% =========================================================

\subsection{Diagramas de secuencia}

Los diagramas de secuencia describen, en orden cronológico, los mensajes que intercambian los participantes de un caso de uso. Se modelaron los ocho flujos que involucran a más partes de la arquitectura: al menos uno por cada módulo de la sección \ref{sec:casos_uso}, incluyendo los que dependen de los proveedores externos seleccionados en la sección \ref{sec:factibilidad_tecnica}.

Los participantes corresponden a los contenedores del diagrama de la figura \ref{fig:diagrama_contenedor}: la aplicación móvil, el servicio de autenticación, la API REST y la base de datos, junto con los sistemas externos. Las cajas azules representan contenedores del sistema y las grises, sistemas externos; debajo de cada una, su línea de vida es la línea discontinua vertical. Las flechas continuas son llamadas, las discontinuas son respuestas y los mensajes están numerados en el orden en que ocurren. Los marcos indican fragmentos combinados: \textbf{alt} para alternativas excluyentes, separadas por una línea discontinua, \textbf{opt} para pasos que solo ocurren si se cumple la condición entre corchetes y \textbf{loop} para pasos que se repiten.

% Los diagramas están dibujados en TikZ (no son imágenes): un archivo por
% flujo en figures/capitulo6_diseno_sistema/6.3_Diseno_estructural_y_de_comportamiento/sd_*.tex.
% Estilos y macros en setup/settings.tex, sección G3.

\subsubsection{Registro y verificación de la cuenta}

La figura \ref{fig:sd_registro} muestra el registro con correo electrónico (CU-01 y CU-02). Después de capturar sus datos, el usuario acepta los términos y el aviso de privacidad en el mismo formulario (RF-12): puede abrirlos y aceptarlos desde la ventana que los muestra, o solo marcar la casilla de aceptación. Al presionar «Registrarse», el consentimiento se envía junto con los datos de la cuenta. El servicio de autenticación rechaza los correos ya registrados; en caso contrario, crea la cuenta sin verificar y envía, mediante Resend, un enlace de un solo uso con vigencia de 10 minutos (RN-14), mientras la aplicación muestra el aviso de verificación pendiente (RF-04).

La figura \ref{fig:sd_verificacion} continúa este flujo con la verificación del correo (CU-03). Si el usuario abre el enlace dentro de su vigencia, la cuenta queda verificada y la aplicación habilita sus funciones (RN-17); si el enlace expiró o ya se usó, la aplicación ofrece reenviarlo (RF-03).

\begin{figure}[htbp]
    \centering
    % =========================================================
% Diagrama de secuencia: Registro de cuenta con correo  (6.3.1)
% CU-01, CU-02 · RF-12, RF-01, RF-02, RF-04 · RN-14
% ---------------------------------------------------------
% Estilos y macros: setup/settings.tex, sección G3.
% La verificación del correo (CU-03) está en sd_verificacion_correo.tex.
% Los textos se parten con \\ para que no rebasen el espacio entre dos
% líneas de vida; [n] reserva la altura de n renglones.
% =========================================================
\begin{tikzpicture}

  % ---------------- Participantes ----------------
  \sdActor{usr}{0}{Usuario}
  \sdContenedor{app}{2.9}{Aplicación móvil}{Flutter}
  \sdContenedor{auth}{6.0}{Autenticación}{Supabase Auth}
  \sdContenedor{bd}{9.1}{Base de datos}{PostgreSQL}
  \sdExterno{mail}{12.2}{Servicio de email}{Resend}
  \sdIniciar

  % ---------------- Registro (CU-02) ----------------
  \sdMensaje[3]{usr}{app}{captura nombre,\\apellidos, correo\\y contraseña}

  % ---------------- Aviso de privacidad (CU-01) ----------------
  % Se acepta desde el mismo formulario: leyéndolo o solo marcando la casilla
  \sdAbrir{aviso}{alt}{[desea leer los términos y el aviso]}{usr}{auth}
    \sdMensaje[2]{usr}{app}{abre «los términos y\\el aviso de privacidad»}
    \sdMensaje[2]{app}{usr}{muestra los términos\\y el aviso}
    \sdMensaje[2]{usr}{app}{presiona «Aceptar\\y continuar»}
  \sdSino{aviso}{[no los lee]}
    \sdMensaje[2]{usr}{app}{marca la casilla\\de aceptación}
  \sdCerrar{aviso}

  \sdMensaje{usr}{app}{presiona «Registrarse»}
  \sdMensaje[3]{app}{auth}{solicita crear la\\cuenta con los datos\\y el consentimiento}
  \sdMensaje{auth}{bd}{busca el correo}
  \sdRespuesta{bd}{auth}{resultado}

  \sdAbrir{correo}{alt}{[correo ya registrado]}{usr}{mail}
    \sdRespuesta[2]{auth}{app}{error:\\correo en uso}
    \sdMensaje[2]{app}{usr}{informa que el\\correo ya existe}
  \sdSino{correo}{[correo disponible]}
    \sdMensaje[3]{auth}{bd}{crea la cuenta no\\verificada y guarda\\el consentimiento}
    \sdMensaje[2]{auth}{mail}{envía el enlace de verificación\\con vigencia de 10 min}
    \sdMensaje{mail}{usr}{entrega el correo con el enlace}
    \sdRespuesta[2]{auth}{app}{cuenta creada\\sin verificar}
    \sdMensaje[2]{app}{usr}{aviso de verificación\\pendiente}
  \sdCerrar{correo}

  \sdTerminar
\end{tikzpicture}

    \caption{Diagrama de secuencia del registro de una cuenta con correo electrónico}
    \label{fig:sd_registro}
\end{figure}

\begin{figure}[htbp]
    \centering
    % =========================================================
% Diagrama de secuencia: Verificación del correo electrónico  (6.3.1)
% CU-03 · RF-02, RF-03 · RN-14, RN-17
% ---------------------------------------------------------
% Estilos y macros: setup/settings.tex, sección G3.
% Continúa el registro de sd_registro_verificacion.tex: el usuario ya
% recibió el correo con el enlace.
% =========================================================
\begin{tikzpicture}

  % ---------------- Participantes ----------------
  \sdActor{usr}{0}{Usuario}
  \sdContenedor{app}{3.2}{Aplicación móvil}{Flutter}
  \sdContenedor{auth}{6.8}{Autenticación}{Supabase Auth}
  \sdContenedor{bd}{10.4}{Base de datos}{PostgreSQL}
  \sdIniciar

  \sdMensaje[2]{usr}{auth}{abre el enlace de verificación\\desde su correo}
  \sdAbrir{enlace}{alt}{[enlace vigente y sin usar]}{usr}{bd}
    \sdMensaje[2]{auth}{bd}{marca el correo\\como verificado}
    \sdRespuesta[2]{auth}{app}{redirige a la\\aplicación}
    \sdMensaje[2]{app}{usr}{habilita las\\funciones}
  \sdSino{enlace}{[enlace expirado o ya usado]}
    \sdRespuesta{auth}{app}{enlace inválido}
    \sdMensaje[2]{app}{usr}{ofrece reenviar\\el enlace}
  \sdCerrar{enlace}

  \sdTerminar
\end{tikzpicture}

    \caption{Diagrama de secuencia de la verificación del correo electrónico}
    \label{fig:sd_verificacion}
\end{figure}

\subsubsection{Registro e inicio de sesión con Google}

En la figura \ref{fig:sd_google} se presenta el acceso con una cuenta de Google (CU-06), disponible tanto en la pantalla de registro como en la de inicio de sesión. El usuario solo confirma su cuenta de Google: no se le pide verificar el correo (RN-17) ni se le muestra el aviso de privacidad. La aplicación obtiene de Google un token de identidad y lo entrega al servicio de autenticación, que verifica su firma, su emisor y su audiencia antes de confiar en él. Si el correo ya pertenece a una cuenta con contraseña, ambas formas de acceso se vinculan a la misma cuenta (RF-06); si es un usuario nuevo, su cuenta se crea en ese momento.

\begin{figure}[htbp]
    \centering
    % =========================================================
% Diagrama de secuencia: Inicio de sesión con Google  (6.3.1)
% CU-06 (incluye CU-01) · RF-06, RF-12 · RN-17
% ---------------------------------------------------------
% Estilos y macros: setup/settings.tex, sección G3.
% =========================================================
\begin{tikzpicture}

  % ---------------- Participantes ----------------
  \sdActor{usr}{0}{Peatón\\usuario}
  \sdContenedor{app}{2.9}{Aplicación móvil}{Flutter}
  \sdExterno{goo}{6.0}{Identidad de Google}{Google Sign-In}
  \sdContenedor{auth}{9.1}{Autenticación}{Supabase Auth}
  \sdContenedor{bd}{12.2}{Base de datos}{PostgreSQL}
  \sdIniciar

  % ---------------- Aviso de privacidad (CU-01) ----------------
  \sdAbrir{aviso}{opt}{[primer uso de la aplicación]}{usr}{bd}
    \sdMensaje[2]{app}{usr}{muestra el aviso\\de privacidad}
    \sdMensaje{usr}{app}{acepta el aviso}
  \sdCerrar{aviso}

  % ---------------- Autorización en Google ----------------
  \sdMensaje[2]{usr}{app}{elige «Continuar\\con Google»}
  \sdMensaje[2]{app}{goo}{solicita el inicio\\de sesión}
  \sdMensaje{goo}{usr}{muestra el selector de cuentas}
  \sdMensaje{usr}{goo}{elige su cuenta y autoriza}
  \sdRespuesta[2]{goo}{app}{token de identidad\\(ID token)}

  % ---------------- Sesión en Supabase ----------------
  \sdMensaje{app}{auth}{signInWithIdToken(token)}
  \sdMensaje[2]{auth}{goo}{valida la firma, el emisor\\y la audiencia del token}
  \sdRespuesta[2]{goo}{auth}{token válido\\(correo y nombre)}
  \sdMensaje[2]{auth}{bd}{busca la cuenta\\por correo}
  \sdRespuesta{bd}{auth}{resultado}

  \sdAbrir{cuenta}{alt}{[el correo ya tiene cuenta con contraseña]}{app}{bd}
    \sdMensaje[2]{auth}{bd}{vincula la identidad\\de Google (RF-06)}
  \sdSino{cuenta}{[usuario nuevo]}
    \sdMensaje[2]{auth}{bd}{crea la cuenta sin\\verificación (RN-17)}
  \sdCerrar{cuenta}

  \sdRespuesta[2]{auth}{app}{sesión (tokens de\\acceso y renovación)}
  \sdMensaje[2]{app}{usr}{muestra el mapa\\principal}

  \sdTerminar
\end{tikzpicture}

    \caption{Diagrama de secuencia del registro e inicio de sesión con Google}
    \label{fig:sd_google}
\end{figure}

\subsubsection{Búsqueda de una dirección}

La figura \ref{fig:sd_busqueda} describe la búsqueda de una dirección de origen o destino (CU-21). Para no saturar los servicios, la aplicación espera 300 ms sin cambios en el texto antes de consultar una Edge Function. Esta busca primero en el índice propio, que solo contiene los lugares del área de servicio, y únicamente consulta al servicio de geocodificación cuando no hay coincidencias suficientes, con Photon como respaldo si Geoapify falla. El área se delimita desde la consulta: a ambos proveedores se les envía el rectángulo que envuelve el área de servicio y nunca la posición del usuario (RN-16). Como estos servicios solo admiten rectángulos y el área es un polígono irregular, la Edge Function descarta al final las sugerencias que caen en las esquinas del rectángulo fuera del polígono (RN-01).

\begin{figure}[htbp]
    \centering
    % =========================================================
% Diagrama de secuencia: Búsqueda de dirección de origen y destino  (6.3.1)
% CU-21 · RF-15, RF-17, RF-20 · RN-01, RN-16
% ---------------------------------------------------------
% Estilos y macros: setup/settings.tex, sección G3.
% Búsqueda en dos niveles (factibilidad técnica, 5.5.1): primero el índice
% propio en PostGIS y, solo si no basta, Geoapify, con Photon de respaldo.
% El sesgo usa el área de servicio, no la posición del usuario (RN-16).
% =========================================================
\begin{tikzpicture}

  % ---------------- Participantes ----------------
  \sdActor{usr}{0}{Peatón\\usuario}
  \sdContenedor{app}{2.9}{Aplicación móvil}{Flutter}
  \sdContenedor{api}{6.0}{API REST}{Supabase Edge Function}
  \sdContenedor{bd}{9.1}{Base de datos}{PostgreSQL, PostGIS}
  \sdExterno{geo}{12.2}{Servicio de geocodificación}{Geoapify, Photon}
  \sdIniciar

  \sdAbrir{teclas}{loop}{[por cada cambio de texto de 3 o más caracteres]}{usr}{geo}
    \sdMensaje{usr}{app}{escribe la dirección}
    \sdPropio[2]{app}{espera 300 ms\\sin teclear}
    \sdMensaje[2]{app}{api}{buscar\_lugares\\(texto)}
    \sdMensaje[2]{api}{bd}{busca en el índice\\propio (pg\_trgm)}
    \sdRespuesta{bd}{api}{coincidencias}

    \sdAbrir[0.55]{externo}{opt}{[coincidencias insuficientes]}{api}{geo}
      \sdMensaje[2]{api}{geo}{autocompletado con filter=rect y\\bias=rect del área de servicio}
      \sdAbrir[0.4]{proveedor}{alt}{[Geoapify responde]}{api}{geo}
        \sdRespuesta{geo}{api}{sugerencias}
      \sdSino{proveedor}{[Geoapify falla]}
        \sdMensaje{api}{geo}{misma consulta a Photon}
        \sdRespuesta{geo}{api}{sugerencias}
      \sdCerrar{proveedor}
      \sdMensaje[2]{api}{bd}{descarta las que están\\fuera del área (RN-01)}
      \sdRespuesta{bd}{api}{sugerencias válidas}
    \sdCerrar{externo}

    \sdRespuesta{api}{app}{sugerencias}
    \sdMensaje[2]{app}{usr}{muestra favoritos y\\sugerencias (RF-17)}
  \sdCerrar{teclas}

  \sdMensaje{usr}{app}{elige una sugerencia}
  \sdMensaje[2]{app}{usr}{fija el origen o el\\destino en el mapa}

  \sdTerminar
\end{tikzpicture}

    \caption{Diagrama de secuencia de la búsqueda de una dirección}
    \label{fig:sd_busqueda}
\end{figure}

\subsubsection{Cálculo de rutas}

El cálculo de rutas (CU-25) se muestra en la figura \ref{fig:sd_rutas}. Sin conexión, la aplicación no realiza la consulta y avisa al usuario (RF-26). Con conexión, invoca mediante la API REST una función de la base de datos que valida que el origen y el destino estén dentro del área de servicio (RN-01), aplica los pesos de la franja horaria de la consulta (RN-05) y calcula con pgRouting las rutas segura, rápida y balanceada (RN-02, RN-03). La respuesta puede indicar un punto fuera de cobertura (RF-20), la inexistencia de una ruta peatonal (RF-25) o las rutas encontradas, de las que la aplicación calcula el tiempo estimado (RN-18) y presenta la distancia y el nivel de riesgo (RF-22).

\begin{figure}[htbp]
    \centering
    % =========================================================
% Diagrama de secuencia: Cálculo de rutas  (6.3.1)
% CU-25 · RF-20, RF-21, RF-22, RF-25, RF-26, RF-27 · RN-01, RN-02, RN-03, RN-05, RN-18
% ---------------------------------------------------------
% Estilos y macros: setup/settings.tex, sección G3.
% La lógica de ruteo vive en una función de la base de datos (pgRouting)
% que la aplicación invoca como RPC a través de la API REST.
% =========================================================
\begin{tikzpicture}

  % ---------------- Participantes ----------------
  \sdActor{usr}{0}{Peatón\\usuario}
  \sdContenedor{app}{3.2}{Aplicación móvil}{Flutter}
  \sdContenedor{api}{6.8}{API REST}{Supabase PostgREST}
  \sdContenedor{bd}{11.4}{Base de datos}{PostGIS, pgRouting}
  \sdIniciar

  \sdMensaje[2]{usr}{app}{solicita las rutas\\(origen y destino)}

  \sdAbrir{red}{alt}{[sin conexión a internet]}{usr}{bd}
    \sdMensaje[2]{app}{usr}{avisa que no hay\\conexión (RF-26)}
  \sdSino{red}{[con conexión]}
    \sdMensaje[2]{app}{api}{rpc calcular\_rutas(origen,\\destino, hora de consulta)}
    \sdMensaje{api}{bd}{ejecuta calcular\_rutas()}
    \sdPropioIzq[2]{bd}{valida el origen y el destino\\en el área de servicio (RN-01)}
    \sdPropioIzq[2]{bd}{toma los pesos de la\\franja horaria (RN-05)}
    \sdPropioIzq[3]{bd}{calcula las rutas segura,\\rápida y balanceada\\(RN-02, RN-03)}

    \sdAbrir[0.55]{resultado}{alt}{[origen o destino fuera del área]}{app}{bd}
      \sdRespuesta{bd}{api}{error de cobertura}
      \sdRespuesta{api}{app}{error de cobertura}
      \sdMensaje[2]{app}{usr}{indica la restricción\\de cobertura (RF-20)}
    \sdSino{resultado}{[no existe ruta peatonal]}
      \sdRespuesta{bd}{api}{sin resultados}
      \sdRespuesta{api}{app}{sin resultados}
      \sdMensaje[2]{app}{usr}{avisa que no existe\\ruta (RF-25)}
    \sdSino{resultado}{[rutas encontradas]}
      \sdRespuesta[2]{bd}{api}{rutas con distancia y\\riesgo por tramo}
      \sdRespuesta[2]{api}{app}{rutas (GeoJSON)\\y resumen}
      \sdPropio[2]{app}{calcula el tiempo a 4.5 km/h (RN-18)\\y marca las rutas iguales (RF-21)}
      \sdMensaje[3]{app}{usr}{muestra las rutas con\\distancia, tiempo y\\riesgo (RF-22)}
    \sdCerrar{resultado}
  \sdCerrar{red}

  \sdTerminar
\end{tikzpicture}

    \caption{Diagrama de secuencia del cálculo de rutas}
    \label{fig:sd_rutas}
\end{figure}

\subsubsection{Navegación con recálculo por desvío}

La figura \ref{fig:sd_navegacion} presenta la navegación activa (CU-29). Si la aplicación aún no tiene el permiso de ubicación, lo solicita (RF-29). Durante la navegación, cada nueva posición actualiza el marcador del usuario (RF-30) y, si tiene seguimientos activos, se publica a sus seguidores. Cuando el usuario se aparta más de 50 m de la ruta (RN-18), la aplicación solicita una nueva ruta desde su posición actual, la traza y avisa del recálculo (RF-31), notificando también a los seguidores.

\begin{figure}[htbp]
    \centering
    % =========================================================
% Diagrama de secuencia: Navegación con recálculo por desvío  (6.3.1)
% CU-29 · RF-29, RF-30, RF-31, RF-33 · RN-18
% ---------------------------------------------------------
% Estilos y macros: setup/settings.tex, sección G3.
% =========================================================
\begin{tikzpicture}

  % ---------------- Participantes ----------------
  \sdActor{usr}{0}{Peatón\\usuario}
  \sdContenedor{app}{2.9}{Aplicación móvil}{Flutter}
  \sdExterno{dev}{6.0}{Servicios del dispositivo}{GPS, Android}
  \sdContenedor{api}{9.1}{API REST}{Supabase}
  \sdContenedor{bd}{12.2}{Base de datos}{PostGIS, pgRouting}
  \sdIniciar

  \sdMensaje[2]{usr}{app}{inicia la navegación\\de la ruta elegida}

  % ---------------- Permiso de ubicación (RF-29) ----------------
  \sdAbrir{permiso}{opt}{[sin permiso de ubicación]}{usr}{dev}
    \sdMensaje[2]{app}{dev}{solicita el permiso\\de ubicación}
    \sdMensaje{dev}{usr}{muestra el diálogo de permiso}
    \sdMensaje{usr}{dev}{concede el permiso}
    \sdRespuesta[2]{dev}{app}{permiso\\concedido}
  \sdCerrar{permiso}
  \sdMensaje[2]{app}{dev}{activa la ubicación\\en primer plano}

  % ---------------- Seguimiento de la posición ----------------
  \sdAbrir{nav}{loop}{[mientras la navegación esté en curso]}{usr}{bd}
    \sdMensaje{dev}{app}{nueva posición}
    \sdPropio[2]{app}{actualiza el marcador\\y el avance (RF-30)}
    \sdAbrir[0.55]{segA}{opt}{[hay seguimientos activos]}{app}{api}
      \sdMensaje[2]{app}{api}{publica la posición (Realtime)}
    \sdCerrar{segA}
    \sdAbrir[0.55]{desvio}{opt}{[desvío mayor a 50 m (RN-18)]}{app}{bd}
      \sdMensaje[2]{app}{api}{rpc calcular\_ruta(posición\\actual, destino, criterio)}
      \sdMensaje[2]{api}{bd}{recalcula con\\pgRouting}
      \sdRespuesta{bd}{api}{nueva ruta}
      \sdRespuesta{api}{app}{nueva ruta}
      \sdMensaje[2]{app}{usr}{traza la nueva ruta y\\avisa el recálculo (RF-31)}
      \sdAbrir[0.4]{segB}{opt}{[hay seguimientos activos]}{app}{api}
        \sdMensaje[2]{app}{api}{notifica el desvío a los\\seguidores (Realtime)}
      \sdCerrar{segB}
    \sdCerrar{desvio}
  \sdCerrar{nav}

  \sdMensaje[2]{app}{usr}{indica la llegada\\al destino (RF-33)}

  \sdTerminar
\end{tikzpicture}

    \caption{Diagrama de secuencia de la navegación con recálculo por desvío}
    \label{fig:sd_navegacion}
\end{figure}

\subsubsection{Registro de un reporte de riesgo}

En la figura \ref{fig:sd_reporte} se muestra el registro de un reporte (CU-31). El usuario indica el lugar con su ubicación actual o con un punto del mapa, elige una categoría del catálogo (RN-06) y puede agregar una descripción. Al insertar el reporte, la base de datos comprueba que la cuenta esté verificada, que el punto esté dentro del área de servicio (RN-01) y que el usuario no haya rebasado 10 reportes en 24 horas (RN-12); si alguna condición falla, la aplicación informa del rechazo (RF-37). Los reportes aceptados quedan en estado pendiente (RN-07) hasta que otros usuarios los confirmen o desmientan, proceso que corresponde a CU-33.

\begin{figure}[htbp]
    \centering
    % =========================================================
% Diagrama de secuencia: Registro de un reporte de riesgo  (6.3.1)
% CU-31 · RF-29, RF-36, RF-37, RF-38 · RN-01, RN-06, RN-07, RN-12
% ---------------------------------------------------------
% Estilos y macros: setup/settings.tex, sección G3.
% La validación por confirmaciones de otros usuarios (CU-33, RN-07)
% ocurre después y no forma parte de este flujo.
% =========================================================
\begin{tikzpicture}

  % ---------------- Participantes ----------------
  \sdActor{usr}{0}{Peatón\\usuario}
  \sdContenedor{app}{2.9}{Aplicación móvil}{Flutter}
  \sdExterno{dev}{6.0}{Servicios del dispositivo}{GPS, Android}
  \sdContenedor{api}{9.1}{API REST}{Supabase PostgREST}
  \sdContenedor{bd}{12.2}{Base de datos}{PostgreSQL, PostGIS}
  \sdIniciar

  \sdMensaje[2]{usr}{app}{abre «Reportar\\incidente»}

  \sdAbrir{lugar}{alt}{[usa su ubicación actual]}{usr}{dev}
    \sdMensaje[2]{app}{dev}{solicita la ubicación\\actual}
    \sdRespuesta{dev}{app}{coordenadas}
  \sdSino{lugar}{[elige un punto en el mapa]}
    \sdMensaje[2]{usr}{app}{selecciona el punto\\en el mapa}
  \sdCerrar{lugar}

  \sdMensaje[3]{usr}{app}{elige la categoría\\(RN-06) y escribe una\\descripción opcional}
  \sdMensaje[2]{app}{api}{inserta el reporte (categoría,\\ubicación y descripción)}
  \sdMensaje[2]{api}{bd}{INSERT con políticas\\de seguridad (RLS)}
  \sdPropioIzq[3]{bd}{verifica la cuenta verificada,\\el área de servicio (RN-01)\\y el límite de 10 en 24 h (RN-12)}

  \sdAbrir{valido}{alt}{[límite alcanzado o fuera del área]}{usr}{bd}
    \sdRespuesta{bd}{api}{rechazo}
    \sdRespuesta{api}{app}{error}
    \sdMensaje[2]{app}{usr}{avisa que alcanzó el\\límite (RF-37)}
  \sdSino{valido}{[reporte válido]}
    \sdPropioIzq[2]{bd}{guarda el reporte en\\estado pendiente (RN-07)}
    \sdRespuesta{bd}{api}{reporte registrado}
    \sdRespuesta{api}{app}{confirmación}
    \sdMensaje{app}{usr}{«Reporte enviado»}
  \sdCerrar{valido}

  \sdTerminar
\end{tikzpicture}

    \caption{Diagrama de secuencia del registro de un reporte de riesgo}
    \label{fig:sd_reporte}
\end{figure}

\subsubsection{Compartir y visualizar un seguimiento}

La figura \ref{fig:sd_seguimiento} describe cómo el usuario comparte su ubicación (CU-35) y cómo la consulta otra persona (CU-38). El enlace se crea con la vigencia elegida por el usuario (RN-13) y se envía con las aplicaciones de mensajería del dispositivo. Al abrirlo, la aplicación del seguidor comprueba su vigencia y, si sigue activo, se suscribe a un canal de Supabase Realtime por el que recibe la posición, la ruta, la distancia, el tiempo y el riesgo del trayecto (RF-45). Si el enlace expiró, está pausado o fue eliminado, se informa al seguidor (RNF-10).

\begin{figure}[htbp]
    \centering
    % =========================================================
% Diagrama de secuencia: Compartir y visualizar un seguimiento  (6.3.1)
% CU-35, CU-38 · RF-29, RF-42, RF-45 · RN-13, RNF-10, RNF-18
% ---------------------------------------------------------
% Estilos y macros: setup/settings.tex, sección G3.
% Seis participantes: cabeceras de 2.1 cm. La ubicación en segundo plano
% y el menú de compartir del sistema se muestran como mensajes propios de
% la aplicación para que el diagrama quepa en el ancho de la página.
% =========================================================
\begin{tikzpicture}

  % ---------------- Participantes ----------------
  \sdActor{usr}{0}{Peatón\\usuario}
  \sdContenedor[2.1cm]{app}{2.65}{Aplicación del usuario}{Flutter}
  \sdContenedor[2.1cm]{api}{5.3}{API REST}{Supabase Realtime}
  \sdContenedor[2.1cm]{bd}{7.95}{Base de datos}{PostgreSQL}
  \sdContenedor[2.1cm]{apps}{10.6}{Aplicación del seguidor}{Flutter}
  \sdActor{seg}{13.25}{Seguidor}
  \sdIniciar

  % ---------------- Generar y compartir (CU-35) ----------------
  \sdMensaje[3]{usr}{app}{genera un enlace\\con la vigencia\\elegida (RN-13)}
  \sdMensaje[2]{app}{api}{crea el\\seguimiento}
  \sdMensaje[3]{api}{bd}{guarda el enlace\\y el código\\(RNF-18)}
  \sdRespuesta{bd}{api}{identificador}
  \sdRespuesta[2]{api}{app}{enlace y\\código}
  \sdPropio[2]{app}{abre el menú de compartir\\(aplicaciones de mensajería)}
  \sdMensaje{app}{seg}{envía el enlace por mensajería}

  % ---------------- Abrir el seguimiento (CU-38) ----------------
  \sdMensaje{seg}{apps}{abre el enlace}
  \sdMensaje[2]{apps}{api}{consulta el seguimiento}
  \sdMensaje[2]{api}{bd}{valida la\\vigencia (RN-13)}
  \sdRespuesta{bd}{api}{estado}

  \sdAbrir{vigente}{alt}{[enlace vigente y activo]}{app}{seg}
    \sdRespuesta[2]{api}{apps}{datos del trayecto}
    \sdMensaje[2]{apps}{api}{se suscribe al canal\\de ubicación}
    \sdAbrir[0.55]{vivo}{loop}{[mientras el enlace siga vigente]}{app}{seg}
      \sdPropio[2]{app}{obtiene la posición\\en segundo plano}
      \sdMensaje{app}{api}{publica la posición}
      \sdMensaje[2]{api}{apps}{reenvía la posición (Realtime)}
      \sdMensaje[3]{apps}{seg}{muestra posición,\\ruta, distancia,\\tiempo y riesgo}
    \sdCerrar{vivo}
  \sdSino{vigente}{[expirado, pausado o eliminado]}
    \sdRespuesta[2]{api}{apps}{enlace no válido}
    \sdMensaje[3]{apps}{seg}{informa el estado\\del enlace\\(RNF-10)}
  \sdCerrar{vigente}

  \sdTerminar
\end{tikzpicture}

    \caption{Diagrama de secuencia para compartir y visualizar un seguimiento}
    \label{fig:sd_seguimiento}
\end{figure}

\subsubsection{Alertas de reportes cercanos}

Finalmente, la figura \ref{fig:sd_alertas} muestra las alertas de reportes cercanos (CU-34). Al activarlas, la aplicación solicita los permisos de ubicación en segundo plano y de notificaciones, obtiene el token del dispositivo en Firebase Cloud Messaging y lo registra en el servidor (RF-41). Cuando un reporte pasa a validado (RN-07), un webhook de la base de datos ejecuta una Edge Function que busca a los usuarios que se encuentran a 1.5 km o menos del reporte (RN-18) y les envía la notificación con prioridad alta.

\begin{figure}[htbp]
    \centering
    % =========================================================
% Diagrama de secuencia: Alertas de reportes cercanos  (6.3.1)
% CU-34 · RF-29, RF-41 · RN-07, RN-18
% ---------------------------------------------------------
% Estilos y macros: setup/settings.tex, sección G3.
% Primera parte: el usuario activa las alertas. Segunda parte: cuando un
% reporte pasa a validado, un webhook de la base de datos dispara una
% Edge Function que notifica por FCM a los usuarios cercanos.
% =========================================================
\begin{tikzpicture}

  % ---------------- Participantes ----------------
  \sdActor{usr}{0}{Peatón\\usuario}
  \sdContenedor[2.1cm]{app}{2.6}{Aplicación móvil}{Flutter}
  \sdExterno[2.1cm]{dev}{5.2}{Servicios del dispositivo}{Android}
  \sdContenedor[2.1cm]{api}{7.8}{API REST}{Supabase Edge Functions}
  \sdContenedor[2.1cm]{bd}{10.4}{Base de datos}{PostGIS}
  \sdExterno[2.1cm]{push}{13.0}{Servicio push}{Firebase Cloud Messaging}
  \sdIniciar

  % ---------------- Activación (RF-41) ----------------
  \sdMensaje[3]{usr}{app}{activa las alertas\\de reportes\\cercanos}
  \sdMensaje[3]{app}{dev}{solicita ubicación\\en segundo plano\\y notificaciones}
  \sdMensaje{dev}{usr}{muestra los diálogos}
  \sdMensaje{usr}{dev}{concede los permisos}
  \sdRespuesta[2]{dev}{app}{permisos\\concedidos}
  \sdMensaje{app}{push}{solicita el token del dispositivo}
  \sdRespuesta{push}{app}{token de FCM}
  \sdMensaje[2]{app}{api}{registra el token\\y la preferencia}
  \sdMensaje{api}{bd}{guarda el token}

  \sdAbrir{fondo}{loop}{[con la aplicación en segundo plano]}{app}{bd}
    \sdMensaje{dev}{app}{posición}
    \sdMensaje[3]{app}{api}{actualiza la\\última ubicación\\conocida}
    \sdMensaje{api}{bd}{guarda la ubicación}
  \sdCerrar{fondo}

  % ---------------- Notificación de un reporte validado ----------------
  \sdMensaje[2]{bd}{api}{webhook: reporte\\validado}
  \sdMensaje[2]{api}{bd}{busca usuarios\\a 1.5 km o menos}
  \sdRespuesta[2]{bd}{api}{tokens de los\\dispositivos}
  \sdAbrir{envio}{loop}{[por cada dispositivo cercano]}{usr}{push}
    \sdMensaje[2]{api}{push}{envía la notificación\\con prioridad alta}
    \sdMensaje{push}{dev}{entrega la notificación}
    \sdMensaje[2]{dev}{usr}{muestra la alerta del\\reporte cercano}
  \sdCerrar{envio}

  \sdTerminar
\end{tikzpicture}

    \caption{Diagrama de secuencia de las alertas de reportes cercanos}
    \label{fig:sd_alertas}
\end{figure}

\subsection{Diagrama de clases}
